%% ============================================================
%%  Gen-Knob Tuner FYP Poster  (Portrait, ~65 x 100 cm)
%%  Compile:  pdflatex poster.tex   (two passes)
%% ============================================================
\documentclass[final,t]{beamer}

\mode<presentation>{}

\usepackage[size=custom,width=70,height=132,orientation=portrait,scale=1.0]{beamerposter}
\usepackage[T1]{fontenc}
\usepackage[utf8]{inputenc}
\usepackage{lmodern}
\usepackage{graphicx}
\usepackage{amsmath,amssymb}
\usepackage{booktabs}
\usepackage{tabularx}
\usepackage{array}
\usepackage{multicol}
\usepackage{multirow}
\usepackage{tikz}
\usepackage[dvipsnames]{xcolor}
\usepackage{pifont}
\usepackage{microtype}
\usepackage{ragged2e}
\usepackage{hhline}
\usepackage{colortbl}

\usetikzlibrary{shapes.geometric,arrows.meta,positioning,calc,fit,backgrounds}

% ============================================================
%  Colour palette (UoM inspired)
% ============================================================
\definecolor{MoraMaroon}{HTML}{5A1D1F}
\definecolor{MoraGold}{HTML}{C8A04A}
\definecolor{DeepBlue}{HTML}{1F4E79}
\definecolor{DeepGreen}{HTML}{2E7D32}
\definecolor{AccentRed}{HTML}{B71C1C}
\definecolor{SoftGray}{HTML}{F5F3EE}
\definecolor{PaperBG}{HTML}{FAFAF7}
\definecolor{DarkInk}{HTML}{1C1C1C}
\definecolor{PaleBand}{HTML}{EFE6D4}

% ============================================================
%  Beamer colour / font setup
% ============================================================
\setbeamercolor{background canvas}{bg=PaperBG}
\setbeamercolor{block title}{bg=MoraMaroon,fg=white}
\setbeamercolor{block body}{bg=white,fg=DarkInk}
\setbeamercolor{block alerted title}{bg=MoraGold,fg=MoraMaroon}
\setbeamercolor{block alerted body}{bg=SoftGray,fg=DarkInk}
\setbeamerfont{block title}{size=\large,series=\bfseries}
\setbeamerfont{block alerted title}{size=\large,series=\bfseries}
\setbeamertemplate{navigation symbols}{}
\setbeamertemplate{caption}[numbered]
\setbeamertemplate{itemize items}{\color{MoraMaroon}\raisebox{0.2ex}{\footnotesize\ding{228}}}

% Custom block look: rounded, with left accent bar
\setbeamertemplate{block begin}{
  \par\vskip0.4ex
  \begin{beamercolorbox}[rounded=false,sep=0.45em,ht=2.1em,leftskip=1em]{block title}%
    \usebeamerfont{block title}\insertblocktitle%
  \end{beamercolorbox}%
  \nointerlineskip%
  \begin{beamercolorbox}[rounded=false,sep=0.75em]{block body}%
    \justifying\setlength{\parskip}{0.25em}%
}
\setbeamertemplate{block end}{\end{beamercolorbox}\vskip0.8em}

\setbeamertemplate{block alerted begin}{
  \par\vskip0.4ex
  \begin{beamercolorbox}[rounded=false,sep=0.45em,ht=2.1em,leftskip=1em]{block alerted title}%
    \usebeamerfont{block alerted title}\insertblocktitle%
  \end{beamercolorbox}%
  \nointerlineskip%
  \begin{beamercolorbox}[rounded=false,sep=0.75em]{block alerted body}%
    \justifying\setlength{\parskip}{0.25em}%
}
\setbeamertemplate{block alerted end}{\end{beamercolorbox}\vskip0.8em}

% ============================================================
%  Helper commands
% ============================================================
\newcommand{\tick}{\textcolor{DeepGreen}{\ding{51}}}
\newcommand{\cross}{\textcolor{AccentRed}{\ding{55}}}
\newcommand{\Hi}[1]{\textbf{\color{MoraMaroon}#1}}
\newcommand{\KW}[1]{\textbf{\color{DeepBlue}#1}}
\newcommand{\OK}[1]{\textbf{\color{DeepGreen}#1}}

% ============================================================
%  Column widths (two columns)
% ============================================================
\newlength{\sepwidth}
\newlength{\colwidth}
\setlength{\sepwidth}{0.020\paperwidth}
\setlength{\colwidth}{0.470\paperwidth}
\newcommand{\sepcol}{\begin{column}{\sepwidth}\end{column}}

% ============================================================
%  Custom header (title bar)
% ============================================================
\setbeamertemplate{headline}{
  \begin{beamercolorbox}[wd=\paperwidth]{block title}
    \begin{tikzpicture}[remember picture,overlay]
      % Background maroon bar
      \fill[MoraMaroon] (0,0) rectangle (\paperwidth, -9);
      % Gold accent stripe
      \fill[MoraGold] (0,-9) rectangle (\paperwidth, -9.4);

      % Left logo
      \node[anchor=north west,inner sep=0pt] at (0.8,-0.9)
        {\includegraphics[width=5.5cm]{figures/uom_logo.png}};

      % Right badge
      \node[anchor=north east,text=white,align=right,inner sep=0pt] at (\paperwidth-1,-1.5) {%
        \fontsize{22}{26}\selectfont\bfseries Department of Computer Science \& Engineering\\[2pt]
        \fontsize{18}{22}\selectfont\mdseries University of Moratuwa, Sri Lanka\\[3pt]
        \fontsize{18}{22}\selectfont\itshape Final Year Project  $\cdot$  2025/2026};

      % Title, centered
      \node[anchor=center,text=white,align=center] at (\paperwidth/2,-5.3) {%
        \fontsize{44}{52}\selectfont\bfseries Gen-Knob Tuner for Relational Databases\\[6pt]
        \fontsize{26}{30}\selectfont\mdseries\itshape
        Domain-Adapted Cost Modelling for Cross-Engine, Cross-Hardware Database Knob Tuning};

      % Authors + supervisors
      \node[anchor=south,text=white,align=center] at (\paperwidth/2,-8.7) {%
        \fontsize{19}{22}\selectfont\bfseries
        Madara Mendis  \textbar{}  Madushan Suriyabandara  \textbar{}  Nisith Dissanayake  \textbar{}  Nithika Dias
        \hspace{1.5em}%
        \fontsize{16}{20}\selectfont\mdseries\itshape
        Supervisors: Dr.\ Srinath Perera (External) $\cdot$ Dr.\ Buddhika Karunaratne (Internal)};
    \end{tikzpicture}
    \vspace{9.5cm}
  \end{beamercolorbox}
}

\setbeamertemplate{footline}{
  \begin{beamercolorbox}[wd=\paperwidth,ht=2.8em,center]{}
    \color{white}\fontsize{14}{16}\selectfont
    \tikz[overlay,remember picture]{\fill[MoraMaroon] (current page.south west) rectangle ([yshift=2.8em]current page.south east);}%
    \bfseries Gen-Knob Tuner \textbullet{} University of Moratuwa \textbullet{} FYP 2025/2026
    \hfill
    \mdseries Contact: \texttt{\{madara.22, madushan.22, nisith.22, nithika.22\}@cse.mrt.ac.lk}
    \hfill
  \end{beamercolorbox}
}

% ============================================================
\title{Gen-Knob Tuner}
\author{FYP Team}
\institute{University of Moratuwa}
\date{}
% ============================================================

\begin{document}
\begin{frame}[t,fragile]

\begin{columns}[t]
\sepcol

% =============================================================
% =============== LEFT COLUMN =================================
% =============================================================
\begin{column}{\colwidth}

% --- Abstract (alerted / highlighted) ---
\begin{alertblock}{\ding{74}\ \ Abstract}
\justifying
\small
Database performance depends critically on \Hi{hundreds of configuration knobs}
whose optimal values vary across workloads, engines and hardware.
State-of-the-art LLM tuners such as \KW{E2ETune} achieve $\sim$51\% speedups
but demand large GPU clusters and tens of thousands of costly workload replays per engine.
\textbf{Gen-Knob Tuner} rebuilds this pipeline around a \Hi{domain-adapted surrogate cost model}
that transfers knowledge across \textit{PostgreSQL $\leftrightarrow$ MySQL} and
\textit{32\,GB $\leftrightarrow$ 64\,GB RAM} domains, cutting labelling effort by \OK{$\sim$95\%}
while staying runnable on \KW{commodity CPU hardware}.
We contribute: (i) a representativeness-aware \Hi{HDBSCAN + Gemini-embedding} workload
selection compressing 2{,}279 workloads to 120, (ii) an \Hi{engine-flag LightGBM}
cost model whose \textit{Joint Weighted Training} reaches \textbf{MAPE\,=\,4.81\%} on MySQL with 15.8$\times$ less
data, and (iii) a first \Hi{DANN + SSDA} ablation on 10{,}190 samples quantifying when
adversarial domain invariance helps vs.\ hurts.
\end{alertblock}

% --- Problem & Motivation ---
\begin{block}{1.\ \ Problem \& Motivation}
\justifying
\small
Modern RDBMS expose $\geq$150 tunable knobs ($>\!10^{100}$ configurations). Existing AI-driven tuners inherit three painful costs:

\vspace{0.15em}
\begin{itemize}
\setlength{\itemsep}{0.1em}
  \item \cross\ \textbf{Label cost}: HEBO needs $\approx$100 DB replays \emph{per workload} $\times$ thousands of workloads.
  \item \cross\ \textbf{Hardware cost}: E2ETune fine-tunes Mistral-7B on 8$\times$A100 GPUs --- inaccessible for most labs.
  \item \cross\ \textbf{Portability cost}: PostgreSQL-trained models collapse on MySQL / different RAM (concept shift).
\end{itemize}

\vspace{0.15em}
\Hi{Goal:} \emph{an accurate, explainable, CPU-deployable knob-tuning stack in which a single high-quality
source domain transfers to new engines and servers with minimal labelled data.}
\end{block}

% --- Three-Phase Pipeline ---
\begin{block}{2.\ \ Gen-Knob Tuner Pipeline}
\centering
\includegraphics[width=\textwidth]{figures/pipeline.pdf}\\[0.2em]
\justifying
\small
Each phase is \Hi{data-efficient by design}: clustering reduces replay cost,
the surrogate cost model replaces costly HEBO evaluations, and RAG-augmented prompts
ground the fine-tuned LM in the PostgreSQL manual for trustworthy knob JSON at inference.
\end{block}

% --- Workload Clustering ---
\begin{block}{3.\ \ Representativeness-Aware Workload Selection}
\justifying
\small
From \textbf{10 benchmarks} (JOB, SSB, SSB-Flat, TPC-H, TPC-DS, TPC-C, SmallBank, Twitter, Wikipedia, YCSB)
we start with 2{,}279 synthetic workloads. SQL literals are replaced with \texttt{?} placeholders to isolate
\emph{structural} signal, then embedded via Google \KW{gemini-embedding-001} (768-d)
and clustered with \KW{HDBSCAN} --- a density-based algorithm that tolerates noise and
auto-selects cluster count. 12 representatives per benchmark are drawn cluster-proportionally.

\vspace{0.1em}
\centering
\includegraphics[width=\textwidth]{figures/clustering.pdf}

\vspace{0.2em}
\justifying\small
\textbf{Why it matters:} naively tuning every workload would require $\sim$228k HEBO evaluations.
Our selection collapses this to \OK{12{,}000 labelled points} that drive the cost model,
which then labels the remaining 2{,}154 workloads virtually
--- a \Hi{$\approx$95\% reduction in physical DB replays} (from $\approx$1{,}900h to $\approx$122h of compute).
\end{block}

% --- DANN architecture ---
\begin{block}{4.\ \ Domain-Adversarial Cost Model (DANN)}
\justifying
\small
To generalise across \textit{concept shift} (MySQL vs.\ PostgreSQL) and \textit{covariate shift} (32\,GB vs.\ 64\,GB RAM)
we adopt \KW{Ganin \& Lempitsky's DANN}. A shared feature extractor $G_f$ is regularised by a
\Hi{Gradient Reversal Layer (GRL)} that negates gradients by $-\lambda$, producing a latent $z$
that confuses the \textit{Domain Classifier} $G_d$ while still minimising the regression loss through $G_y$.
Semantic normalisation maps \texttt{shared\_buffers}\,$\rightarrow$\,\texttt{memory\_buffer\_size}
so both engines share one feature schema.

\vspace{0.2em}
\centering
\includegraphics[width=0.95\textwidth]{figures/dann_arch.pdf}
\end{block}

\end{column}
\sepcol


% =============================================================
% =============== RIGHT COLUMN ================================
% =============================================================
\begin{column}{\colwidth}

% --- Experimental setup (benchmarks + knobs) ---
\begin{block}{5.\ \ Experimental Setup}
\justifying\small
\begin{center}
\renewcommand{\arraystretch}{1.1}
\footnotesize
\begin{tabularx}{\textwidth}{@{}l c c c c@{}}
\toprule
\rowcolor{PaleBand}
\textbf{Benchmark} & \textbf{Type} & \textbf{Tables} & \textbf{Queries} & \textbf{Synth.\ workloads} \\
\midrule
JOB (IMDb)       & OLAP & 21 & 113 & 243 \\
SSB / SSB-Flat   & OLAP &  5 &  13 & 239 \\
TPC-H            & OLAP &  8 &  22 & 232 \\
TPC-DS           & OLAP & 24 &  99 & 271 \\
TPC-C, SmallBank, Twitter, Wiki, YCSB & OLTP & -- & -- & $\approx$240 ea. \\
\bottomrule
\end{tabularx}
\end{center}

\textbf{Servers:} Intel Xeon E3-1275v5 (64\,GB) + Intel Core i7-7700 (32\,GB), PostgreSQL 12.2 and MySQL 8.0. \textbf{Knobs:} 15 high-impact PostgreSQL (\texttt{shared\_buffers}, \texttt{work\_mem}, \texttt{effective\_cache\_size}, \ldots) + 10 MySQL analogues. \textbf{Features:} 15 internal metrics + 32-d PCA of EXPLAIN-plan embed + 10 workload descriptors $=$ \Hi{65-d vector}. \textbf{Labels:} 100 \KW{HEBO} iterations per workload.
\end{block}

% --- Regression results ---
\begin{block}{6.\ \ Results --- Cost Model Accuracy}
\textbf{A.\ PostgreSQL base model \ (5-fold \textsc{GroupKFold} on unseen workloads; 217 workloads, 20{,}758 samples)}

\vspace{0.1em}
\centering
\includegraphics[width=0.93\textwidth]{figures/postgres_algos.pdf}

\justifying\small
\textbf{LightGBM} wins across RMSE and MAPE while remaining $10\times$ smaller than an MLP surrogate; its importance ranking also aligns with DBA intuition (\texttt{work\_mem}, \texttt{shared\_buffers}, \texttt{effective\_cache\_size} dominate).

\vspace{0.3em}
\textbf{B.\ Cross-engine transfer  (PostgreSQL $\rightarrow$ MySQL, 1{,}314 MySQL samples)}
\vspace{0.1em}
\centering
\includegraphics[width=0.93\textwidth]{figures/transfer_learning.pdf}

\justifying\small
\vspace{-0.3em}
\begin{itemize}
\setlength{\itemsep}{0.1em}
\item \Hi{MySQL-only}: poor (MAPE 12.3\%, $\rho$=0.47) --- too little data to learn on its own.
\item \KW{Residual Adapter}: freezes PG model, learns only $\Delta$-latency correction. \tick\ Halves MAPE.
\item \OK{Joint Weighted Training}: upweights MySQL rows by $19.96\times$; best overall (MAPE 4.81\%, $\rho$=0.61).
\end{itemize}
\end{block}

% --- DANN ablation results ---
\begin{block}{7.\ \ Results --- DANN Ablation Study}
\justifying\small
We ablate \textbf{normalisation} (Global vs.\ Per-Engine) against \textbf{adaptation objective}
on 10{,}190 samples split across 4 domains (MySQL-32, MySQL-64, PostgreSQL-32, PostgreSQL-64;
MySQL mean cost 16.6, PostgreSQL 2.57 --- a $\sim$6.5$\times$ scale gap).

\vspace{0.3em}
\begin{center}
\renewcommand{\arraystretch}{1.18}
\footnotesize
\begin{tabularx}{\textwidth}{@{}l l c c c@{}}
\toprule
\rowcolor{PaleBand}
\textbf{Norm.} & \textbf{Experiment} & \textbf{RMSE} & \textbf{MAPE (\%)} & \textbf{Spearman} \\
\midrule
\multirow{3}{*}{Global}
  & (A) Source-only DANN      & 46.31 & 108.12 & 0.732 \\
  & (B) SSDA 10\% labels      & \textbf{45.69} & 246.18 & 0.631 \\
  & (C) Supervised baseline   & 45.70 & \textbf{180.59} & 0.779 \\
\midrule
\multirow{3}{*}{Per-Engine}
  & (A) Source-only DANN      & 109.99 & 1041.05 & 0.804 \\
  & (B) SSDA 10\% labels      & 47.06 & 203.71 & 0.726 \\
  & \cellcolor{SoftGray}(C) Supervised $\bigstar$ & \cellcolor{SoftGray}\textbf{46.37} & \cellcolor{SoftGray}\textbf{99.90} & \cellcolor{SoftGray}\textbf{0.843} \\
\bottomrule
\end{tabularx}
\end{center}

\vspace{0.1em}
\centering
\includegraphics[width=0.93\textwidth]{figures/dann_ablation.pdf}

\justifying\small
\textbf{Key findings:}
\begin{itemize}
\setlength{\itemsep}{0.1em}
\item \tick\ \Hi{Per-engine normalisation + supervised training} gives the strongest ranking ($\rho$\,=\,0.843) --- critical for Top-K config selection.
\item \tick\ Global normalisation favours \KW{SSDA} when engines share scale; \textbf{10\% labelled MySQL data} beats source-only DANN on every run.
\item \cross\ Source-only DANN is \textit{unstable} under per-engine norm (MAPE $>$ 1000\%) --- adversarial invariance alone cannot bridge a 6.5$\times$ scale gap.
\item \tick\ Higher domain-classifier accuracy does \emph{not} imply better regression --- \textit{decision} we avoid over-regularising $G_f$.
\end{itemize}
\end{block}

% --- Contributions & Conclusion combined ---
\begin{alertblock}{8.\ \ Key Contributions \& Take-aways}
\small
\justifying
\begin{itemize}
\setlength{\itemsep}{0.2em}
  \item \tick\ \Hi{Representativeness-aware labelling} (HDBSCAN + Gemini): Silhouette \textbf{0.828}, \OK{$-$95\% HEBO cost}.
  \item \tick\ \Hi{Engine-flag LightGBM} cost model: $\rho$\,=\,\textbf{0.942} on PostgreSQL; \textbf{Joint Weighted Training} achieves first sub-5\% MAPE cross-engine transfer (\textbf{4.81\%}).
  \item \tick\ First systematic \Hi{DANN + SSDA ablation} on a DB-cost-model task --- per-engine norm $\bigstar$ wins ($\rho$\,=\,0.843).
  \item \tick\ End-to-end runnable on \Hi{a single CPU} --- knob JSON in $<$5\,s via a quantised 1.5B SLM + RAG over PG docs.
\end{itemize}

\vspace{0.3em}
\textbf{Take-away:}\ Adversarial DANN is best when the target engine has \emph{zero} labels and scales are close;
otherwise \Hi{per-engine normalisation + a small weighted sample} of the target engine matches or beats it at a fraction of the complexity.
\end{alertblock}

% --- Future work + References (compact) ---
\begin{block}{9.\ \ Future Work \ \ \& \ \ Selected References}
\small
\justifying
\textbf{Short term:}\ Full SFT of Qwen-2.5 (7B) / DeepSeek-R1 (1.5B) on cost-model-guided labels; RAG validation of generated knob JSON. \quad
\textbf{Long term:}\ Real-time dynamic tuning and integration with learned indexing / cardinality estimation; adaptation to distributed DBs (CockroachDB, TiDB).

\vspace{0.3em}
\fontsize{11}{13}\selectfont
\textbf{[1]}\ X.\ Huang \emph{et al.}, E2ETune, \emph{arXiv:2404.11581}, 2024. \ 
\textbf{[2]}\ Y.\ Ganin \& V.\ Lempitsky, Unsupervised Domain Adaptation by Backprop, \emph{ICML}, 2015. \
\textbf{[3]}\ A.\ Cowen-Rivers \emph{et al.}, HEBO, \emph{JAIR}, 2022. \
\textbf{[4]}\ D.\ Van Aken \emph{et al.}, OtterTune, \emph{SIGMOD}, 2017. \
\textbf{[5]}\ R.\ Campello \emph{et al.}, HDBSCAN, \emph{PAKDD}, 2013. \
\textbf{[6]}\ G.\ Ke \emph{et al.}, LightGBM, \emph{NeurIPS}, 2017.
\end{block}

\end{column}
\sepcol
\end{columns}

\end{frame}
\end{document}
